% Chapter 6. Historical FCP BA-BH source native operator chain.
\chapter[Native Physical Hamiltonian]{The Native Physical Closed Charge and Self-Adjoint Hamiltonian}
\label{ch:hamiltonian}

\section{The problem left by the smooth-core equations}
The charge formula \eqref{eq:source-charge} is initially defined on a
smooth test-function core. The physical Hilbert space of
Chapter~\ref{ch:gauge} is an $L^2$ closed subspace. To speak rigorously
about the spectrum of a physical Hamiltonian, we need an actual densely
defined closed operator on that physical Hilbert space, with its
operator domain specified. A formal integration-by-parts calculation on
$C_c^\infty$ is insufficient: distinct closed extensions of an
unbounded differential expression can have different spectra.

The central achievement of the October FCP physical-domain program was to
construct this operator \emph{from the original restricted charge graph}
and to compare the resulting native Hamiltonian with a historically
constructed paired one. This was not merely an isolated proposition; it
required a sequence of exact graph, gauge-core, density, adjoint, and
natural-domain results.

\section{The native physical charge map}
Let $\mathscr H_{\mathrm{phys}}\subset\mathscr H$ be as above and take
$\mathscr D_{\mathrm{phys}}\subset\mathscr D$ to be the source
invariant core. Restrict the original sixteen-charge column
$\mathscr Q$ to this core. The historical source constructs its \emph{minimal
closed} graph in the relevant product of the actual physical input Hilbert
space and the sixteen-component output Hilbert space. Denote the
resulting operator by
\begin{equation}\label{eq:physical-T}
 T:\operatorname{Dom}(T)\subset\mathscr H_{\mathrm{phys}}
 \longrightarrow \mathscr K,
 \qquad \mathscr K:=\bigoplus_{\alpha=0}^{15}\mathscr H.
\end{equation}
We do not replace its graph by the closure of the unrestricted source
\texttt{fullClosedGraph}, nor assert that $T$ is a bounded map defined on
all of $\mathscr H_{\mathrm{phys}}$.

On $\mathscr D_{\mathrm{phys}}$ the source identifies
$T\iota\psi=(\iota Q_\alpha\psi)_\alpha$.
The closed operator is the \emph{minimal} extension obtained by graph
closure. Its underlying source declarations are the AY original physical
Hilbert charge and the earlier AU closed graph. Their formalization
records are indexed in the companion theorem atlas.

\section{Density is a theorem, not a convenience}
\begin{sourceaccount}[Original physical charge domain density;
historically qualified]\label{thm:T-dense}
For the original $N=2$ color gauge system, the natural domain of the
actual original minimal closed physical Hilbert charge column is dense:
\begin{equation}\label{eq:T-dense}
 \overline{\operatorname{Dom}(T)}
 =\mathscr H_{\mathrm{phys}}.
\end{equation}
\end{sourceaccount}
\begin{sourceexplanation}
The original AY source operator is the graph closure of the literal
restricted invariant-core charge. An earlier AZ reduction expresses its
domain density as a consequence of an explicit equality between the
source invariant-core norm closure and the whole physical subspace.
The historical D2B10C theorem independently establishes precisely that
core-norm equality using the original SU(2) Haar construction. BA
composes the two results in one pinned source environment, thereby
eliminating the formerly conditional density premise. It does not
replace the physical domain with a convenient larger test space.
\end{sourceexplanation}
\formalrecord{Original physical charge domain density theorem BA;
full Lean name in atlas H-01. FCP strict
\href{https://github.com/etblink/Foundational-Convergence-Program/actions/runs/39079626713}{\#39079626713}
(17 October 2026).}

This is a pivotal improvement over earlier records: a contemporaneous
October 17 AW/AX handoff still lists density as an \emph{open} task.
That handoff is a faithful snapshot of the research \emph{at that time},
but it is not the latest theorem status. A narrative reconstructed solely
from the earlier note would mistakenly describe an already-completed
mathematical obligation as open.

\section{The adjoint square has a particular domain}
For a densely defined closed map $T:\mathscr H_{\mathrm{phys}}\to\mathscr K$,
its Hilbert adjoint is the actual operator
$T^\dagger:\operatorname{Dom}(T^\dagger)\subset\mathscr K
\to\mathscr H_{\mathrm{phys}}$. The natural composition is
\begin{equation}\label{eq:TstarT-domain}
 \operatorname{Dom}(T^\dagger T)
 =\{\psi\in\operatorname{Dom}(T):
 T\psi\in\operatorname{Dom}(T^\dagger)\},\qquad
 T^\dagger T\psi=T^\dagger(T\psi).
\end{equation}
The second condition is essential. It would be false to infer from
$\psi\in\operatorname{Dom}(T)$ that $T\psi$ automatically belongs to
$\operatorname{Dom}(T^\dagger)$.

\begin{sourceaccount}[Native normalized Hamiltonian;
historically qualified]\label{thm:source-H}
For the actual original physical source operator $T$,
\begin{equation}\label{eq:source-H}
 H_{\mathrm{phys}}:=\frac1{16}T^\dagger T,
 \qquad
 \operatorname{Dom}(H_{\mathrm{phys}})=
 \operatorname{Dom}(T^\dagger T)
\end{equation}
is self-adjoint, nonnegative, densely defined and closed on
$\mathscr H_{\mathrm{phys}}$. Its action and normalized weak-form
representing-vector criterion are identified with the literal original
source definitions.
\end{sourceaccount}
\begin{sourceexplanation}
The historical BA theorem supplies the density needed to form the
Hilbert adjoint. The original minimal closed charge construction supplies
closedness. The standard adjoint-square theorem for a densely defined
closed Hilbert-space operator gives that $T^\dagger T$ is positive and
self-adjoint on the natural domain \eqref{eq:TstarT-domain}; in the FCP
formalization, BC constructs that product and BD proves self-adjointness.
Multiplication by the positive real scalar $1/16$ preserves that domain
and self-adjointness. BE proves the exact normalized scalar action,
closedness, density, and characterization by the weak form. This is a
precise source-typed application of the general theorem, not an
assumption that a formal differential Hamiltonian happens to be
essentially self-adjoint on test functions.
\end{sourceexplanation}
\formalrecord{Original native normalized Hamiltonian theorem family
BE (self-adjointness, natural domain and weak form); exact declarations
in atlas H-02. FCP strict
\href{https://github.com/etblink/Foundational-Convergence-Program/actions/runs/39088799248}{\#39088799248}.}

\section{The weak form fixes the associated operator}
The actual closed physical quadratic form is
\begin{equation}\label{eq:closed-q}
 q(\psi):=\frac1{16}\|T\psi\|_{\mathscr K}^2,
 \qquad \psi\in\operatorname{Dom}(T),
\end{equation}
and its polarized form is
\begin{equation}\label{eq:polarized-q}
 q(\psi,\phi):=\frac1{16}
 \langle T\psi,T\phi\rangle_{\mathscr K},\qquad
 \psi,\phi\in\operatorname{Dom}(T).
\end{equation}
There is a precise test for whether a vector $\psi$ lies in the
domain of the associated operator: it must admit an actual
$w\in\mathscr H_{\mathrm{phys}}$ with
\begin{equation}\label{eq:weak-form-criterion}
 \langle w,\phi\rangle_{\mathscr H_{\mathrm{phys}}}
 =\frac1{16}\langle T\psi,T\phi\rangle_{\mathscr K}
 \quad\text{for every }\phi\in\operatorname{Dom}(T).
\end{equation}
By the Hilbert adjoint definition and density of $\operatorname{Dom}(T)$,
this is equivalent to $T\psi\in\operatorname{Dom}(T^\dagger)$, in which
case $w=\frac1{16}T^\dagger T\psi$. Thus the weak form identifies the
\emph{same natural-domain operator}, not an arbitrary extension.

The earlier D2B23--D2B25 branch also gave a concrete complete weighted
Hilbert graph representation of this source form, including the exact
form norm. The graph norm equivalence and weighted-closure results are
useful in later compactness arguments, but they must not be confused
with equality between arbitrary operator domains.
\formalrecord{FCP strict weighted-form stages
\href{https://github.com/etblink/Foundational-Convergence-Program/actions/runs/38060943895}{\#38060943895},
\href{https://github.com/etblink/Foundational-Convergence-Program/actions/runs/38061745424}{\#38061745424}, and
\href{https://github.com/etblink/Foundational-Convergence-Program/actions/runs/38062563927}{\#38062563927}.
The later BE weak-form theorem supplies the source-normalized native
operator.}

\section{The historical and native Hamiltonians were identified}
The physical model existed in two separate source-defined lines: the
original AY/BE minimal physical Hilbert charge construction, and the
earlier historical D2B17--D2B28 paired construction. It was
mathematically impermissible merely to point to their similar formulas
and declare their closed operators equal. The FCP BF and BG stages first
bound the charge graphs and domains. BH then compared adjoints,
compositions and normalized operators.

\begin{sourceaccount}[Original and paired Hamiltonian identity;
historically qualified]\label{thm:BH}
At the historically accepted original source snapshot, the AY/BE native
and D2B historical charge adjoints, natural adjoint squares, and
normalized Hamiltonians agree as the exact source-defined unbounded
operators, including their natural domains. On the actual paired SU(2)
source data, the two normalized physical Hamiltonian declarations agree.
\end{sourceaccount}
\begin{sourceexplanation}
BG identifies the genuine original and historical physical charge
operators as \emph{partial linear maps}, including their domains and
closed graphs. Their Hilbert adjoints consequently agree, with the
same adjoint domains. BH rewrites the two separately defined
natural compositions into those identified adjoints and charge
operators; equality of the compositions gives equality of their
domains. Multiplying by the common scalar $1/16$ then establishes
equality of the normalized Hamiltonians, not simply coincidence of
their test-function formulas. The final paired specialization uses
$\texttt{pairedAlgebraData}$ and $\texttt{pairedGaugeData}$ rather than
a substitute representation.
\end{sourceexplanation}
\formalrecord{Source-identification theorem family BH (adjoints,
natural squares and Hamiltonian equality with domains); exact seven
Lean declaration names in atlas H-03. Historical strict
\href{https://github.com/etblink/Foundational-Convergence-Program/actions/runs/39096180117}{\#39096180117}
(17 October 2026).}

\section{A crucial porting distinction: accepted source versus current tree}
The historically qualified BH source used a local \texttt{letI} complete
Hilbert instance. In the present imported \texttt{OperatorDependencies}
copy, the analogous proof line is \texttt{let}, alongside rewritten
imports. This is \emph{not} a demonstrated mathematical counterexample:
the source-level BH theorem has its accepted exact-run certificate.
The separate M3H import history shows how an identical sort of apparent
difference can arise from a strict style-linter warning. But neither
historical success nor the M3H experience is a substitute for a fresh
independent current-port consumer test. We therefore classify BH as
\emph{historically strict-qualified; current modified import requires its
own exact qualification}.

\section{What we know---and what we do not}
We have a genuine source physical Hilbert space, a densely defined
closed original charge column, its actual adjoint, a positive
self-adjoint normalized Hamiltonian with precisely the natural domain,
and a historically qualified identity to the paired construction.
That settles an important part of the operator-theory foundation.
It does \emph{not} yet construct a continuous spatial $\operatorname{Spin}(9)$
unitary action that preserves this same closed physical form, select
nontrivial high-type reducing sectors, derive radial coercivity in
those sectors, or establish eigenpairs. These are the later chapters'
real proof obligations. Treating them as already solved would erase
the very mathematical frontier this monograph is intended to expose.
